\section{Is It Dynamic Programming? Case Studies}

\begin{frame}{How this part works}
	For each problem you will see \alert{candidate solutions}. Your job is to decide:
	\begin{enumerate}
		\item Is it \textbf{correct}?
		\item Is it \textbf{the right tool}?
	\end{enumerate}

	\seprule

	\textbf{Triage test for dynamic programming}
	\begin{enumerate}
		\item Is there a \emph{sufficient} state of polynomial size? (Markov property)
		\item Does the objective decompose over it? (optimal substructure)
		\item Is there no cheaper algorithm that exploits more structure?
	\end{enumerate}
	A positive answer to (1) and (2) makes DP \emph{correct}; (3) decides whether it is \emph{good}.
\end{frame}

\begin{frame}{Case 1: Split Array Largest Sum (LeetCode 410)}
	Given $a_1,\dots,a_n\ge 0$ and $k\le n$, split the array into $k$ non-empty
	contiguous parts so as to \alert{minimise the largest part sum}. Let $S_i=\sum_{l\le i}a_l$.

	\seprule

	\textbf{Example.} $a=(7,2,5,10,8)$, $k=2$: the parts $(7,2,5)\,|\,(10,8)$ give $\max\{14,18\}=18$, which is optimal.

	\vspace{1em}
	We will look at two proposals.
\end{frame}

\begin{frame}{Case 1: candidate solutions}
	\begin{columns}[T,onlytextwidth]
		\column{0.52\textwidth}
		\begin{block}{Proposal A (recurrence)}
			\small
			Let $f(i,j)$ be the best value for the first $i$ elements in $j$ parts.
			\begin{align*}
				f(0,0)&=0, \quad f(i,0)=f(0,j)=\infty\\
				f(i,j)&=\min_{j-1\le p<i}\max\bigl\{f(p,j-1),\;S_i-S_p\bigr\}
			\end{align*}
			Answer $f(n,k)$.
		\end{block}

		\column{0.44\textwidth}
		\begin{block}{Proposal B (greedy)}
			\small
			Let $\bar S=\lceil S_n/k\rceil$. Sweep left to right and close the current part
			as soon as adding the next element would exceed $\bar S$ (until $k-1$ parts are closed).
		\end{block}
	\end{columns}

	\audience{Vote}{%
		For each proposal: \textbf{correct} or \textbf{incorrect}? And if correct, is it \textbf{efficient}?
	}
\end{frame}

\begin{frame}{Case 1: verdict}
	\begin{columns}[T,onlytextwidth]
		\column{0.5\textwidth}
		\textbf{A is correct.} The state $(i,j)$ is sufficient (only the prefix length and number of parts matter),
		and $(\max,\min)$ composes. But the cost:
		\[
			\underbrace{nk}_{\text{states}}\times\underbrace{O(n)}_{\text{per state}}=\Theta(kn^2).
		\]

		\column{0.46\textwidth}
		\textbf{B is wrong.} Take $a=(1,3,2)$, $k=2$: $\bar S=3$, so B closes $(1)$ and returns $(1)\,|\,(3,2)$ with
		value $5$, while $(1,3)\,|\,(2)$ achieves $4$.
	\end{columns}

	\seprule

	Greedy heuristics that ``balance'' parts fail because the objective is the \emph{maximum}, not a
	deviation from the mean. Is there a correct greedy at all?
\end{frame}

\begin{frame}{Case 1: a better algorithm, and why DP is wasteful}
	\small
	\textbf{Decide, then optimise.} For a cap $L$ let $c(L)$ be the minimum number of parts
	with every sum $\le L$. The greedy ``extend each part as far as possible'' computes $c(L)$ in $O(n)$
	(it stays ahead of any other partition).

	\begin{itemize}
		\item $c$ is non-increasing in $L$; since $k\le n$, a partition into fewer parts can always be refined.
		\item The answer is $\min\{L\in[\max_i a_i,\,S_n]\;:\;c(L)\le k\}$, found by binary search.
	\end{itemize}
	\vspace{0.3em}
	\begin{center}
		\begin{tabular}{@{}lc@{}}
			\toprule
			Method & Time\\
			\midrule
			Proposal A & $\Theta(kn^2)$\\
			A with monotone optimum (binary search over $p$) & $O(kn\log n)$\\
			Binary search on the answer + greedy & $O(n\log S_n)$\\
			\bottomrule
		\end{tabular}
	\end{center}
	DP computes $f(i,j)$ for \emph{every} prefix and part count, but we need one number.
\end{frame}

\begin{frame}{Case 1: discussion}
	\audience{Discuss}{%
		\begin{enumerate}
			\item Which property of ``minimise the \emph{maximum}'' made the decision version
			greedy-checkable and monotone in $L$?
			\item Suppose we instead minimise the \textbf{sum of squares} of part sums. Does binary search
			on the answer still work? What would you use?
		\end{enumerate}
	}

	\uncover<2->{%
		\small
		\begin{enumerate}
			\item The max commutes with refinement: splitting a part never increases the objective,
			and feasibility under a cap is local.
			\item No: there is no threshold structure. This is a genuine DP,
			$f(i,j)=\min_p f(p,j-1)+(S_i-S_p)^2$, and its transition has the form that the
			\alert{convex hull trick} exploits in Part~3.
		\end{enumerate}
	}
\end{frame}

\begin{frame}{Case 2: Maximum weight independent set on a tree}
	LeetCode 337 (\emph{House Robber III}); the path case is LeetCode 198.

	Tree $T=(V,E)$ rooted at $r$, weights $w_v\ge0$. Choose an independent set $I$ maximising $\sum_{v\in I}w_v$.
	Write $\ch(v)$ for the children and $\gc(v)$ for the grandchildren of $v$.

	\seprule

	\begin{itemize}
		\item For general graphs the problem is NP-hard.
		\item On a tree, the subtree of $v$ is separated from the rest of the tree by the single vertex $v$.
	\end{itemize}

	\begin{center}
		Where would you put the state?
	\end{center}
\end{frame}

\begin{frame}{Case 2: candidate solutions}
	\small
	\begin{description}
		\item[\alert{Proposal A}] Greedy: repeatedly take the heaviest remaining vertex and delete its neighbours.
		\item[\alert{Proposal B}] $f(v)=\max\bigl\{\,w_v+\sum_{c\in\ch(v)}f(c),\;\sum_{c\in\ch(v)}f(c)\bigr\}$
		\item[\alert{Proposal C}] $f(v)=\max\bigl\{\,w_v+\sum_{g\in\gc(v)}f(g),\;\sum_{c\in\ch(v)}f(c)\bigr\}$
		\item[\alert{Proposal D}] $f(v,1)=w_v+\sum_{c}f(c,0)$ and $f(v,0)=\sum_c\max\{f(c,0),f(c,1)\}$; answer $\max_b f(r,b)$.
	\end{description}

	\audience{Vote}{%
		Which proposals are \textbf{correct}? Look for a counterexample for those you reject.
	}
\end{frame}

\begin{frame}{Case 2: verdict}
	\small
	\begin{columns}[T,onlytextwidth]
		\column{0.5\textwidth}
		\textbf{A and B are wrong.} Path $2-3-2$:
		\begin{itemize}
			\item A takes the centre: value $3$; optimum $4$ (both ends).
			\item B, rooted at the centre: $f(\text{leaf})=2$, so $f(\text{root})=\max\{3+4,4\}=7$.
			The term $w_v+\sum_c f(c)$ forgets that $f(c)$ may itself use $c$.
		\end{itemize}

		\column{0.46\textwidth}
		\textbf{C and D are correct.} In D the state records whether $v$ is chosen:
		\begin{align*}
			f(v,1)&=w_v+\sum_{c}f(c,0)\\
			f(v,0)&=\sum_{c}\max_b f(c,b)
		\end{align*}
		C is D with the state projected out: $f(v)=\max_b f(v,b)$.
	\end{columns}

	\seprule

	Time $\Theta(\sum_v\deg v)=\Theta(n)$. Both formulations have the same state DAG (the rooted tree), but
	D carries $2n$ states and has \emph{local} transitions. Enlarging the state to make the transition local is a standard trick.
\end{frame}

\begin{frame}{Case 2: why trees, and what generalises}
	The tree makes $v$ a \alert{separator}: the two sides interact only through $v$'s status.

	\begin{theorem}[folklore]
		Given a tree decomposition of width $\omega$ of an $n$-vertex graph, a maximum weight independent
		set can be computed in $O(2^{\omega}\,\omega\, n)$ time.
	\end{theorem}
	The state of a bag is the independent subset chosen inside it ($\le 2^{\omega+1}$ states);
	the structure of the graph \emph{is} the shape of the state DAG.

	\seprule

	\textbf{Same recipe, different shape.}
	Trees ($\omega=1$), cycles ($\omega=2$), series-parallel graphs ($\omega\le2$), grids $k\times n$ ($\omega=k$).
\end{frame}

\begin{frame}{Case 2: discussion}
	\audience{Discuss}{%
		Now let the graph be a \textbf{cycle} $C_n$ with weights $w_1,\dots,w_n$.
		The tree recurrence is no longer valid. Fix it, and give the complexity.
	}

	\uncover<2->{%
		\small
		Fix the status $b_1$ of vertex $1$ and run the path DP on $1,\dots,n$ with that boundary condition
		(for $b_1=1$ vertex $n$ is forbidden): two path DPs, $\Theta(n)$ in total.
		Equivalently, enlarge the state to $(b_1,b_{\text{current}})$: this is the width-$2$ tree decomposition
		made explicit, with $4$ states per position.
	}
\end{frame}

\begin{frame}{Case 3: Minimum Initial Energy to Finish Tasks (LeetCode 1665)}
	Tasks $i=1,\dots,n$ with $(a_i,m_i)$, $a_i\le m_i$: to \emph{start} task $i$ you need at least $m_i$ energy,
	and finishing it consumes $a_i$. Choose the order $\pi$ and the initial energy $E$ to finish all tasks, minimising $E$.

	\seprule

	For an order $\pi$:
	\[
		E(\pi)=\max_{t}\Bigl(m_{\pi(t)}+\sum_{s<t}a_{\pi(s)}\Bigr),\qquad
		E^\star=\min_{\pi}E(\pi).
	\]
	The order is a free choice. There are $n!$ orders.

	\begin{center}
		It looks like a perfect candidate for DP: a sequence of decisions, additive consumption.
	\end{center}
\end{frame}

\begin{frame}{Case 3: candidate solutions}
	Let $A(S)=\sum_{i\in S}a_i$.
	\small
	\begin{description}
		\item[\alert{Proposal A}] Sort by $m_i$ in \emph{decreasing} order and do tasks in that order.
		\item[\alert{Proposal B}] Sort by $a_i$ in \emph{increasing} order.
		\item[\alert{Proposal C}] $F(\emptyset)=0$, $\;F(S)=\min_{i\in S}\max\bigl\{F(S\setminus\{i\}),\;A(S\setminus\{i\})+m_i\bigr\}$; answer $F([n])$.
	\end{description}

	\audience{Vote}{%
		Correct or incorrect? Is C a ``dynamic programming solution''? Is it a \emph{useful} one?
	}
\end{frame}

\begin{frame}{Case 3: verdict}
	\small
	Tasks written $[a_i,m_i]$.
	\begin{itemize}
		\item \textbf{A is wrong:} $[1,2],[3,3]$. A runs $[3,3]$ first and needs $\max\{3,\,3+2\}=5$;
		running $[1,2]$ first needs $\max\{2,\,1+3\}=4$.
		\item \textbf{B is wrong:} $[1,1],[2,3]$. B needs $\max\{1,1+3\}=4$; the other order needs $\max\{3,2+1\}=3$.
		\item \textbf{C is correct.} The energy left after finishing a set $S$ is $E-A(S)$ \emph{whatever the order}, so
		the set is a sufficient state.
	\end{itemize}

	\seprule

	But the state is the \alert{whole subset}: $2^n$ states, $\Theta(2^nn)$ time. No polynomial-size sufficient state
	is apparent: the future depends on \emph{which} tasks remain, not on a number summarising them.
	This is the sense in which the DP is a false positive: the formulation is valid, but the state does not compress.
\end{frame}

\begin{frame}{Case 3: what actually solves it}
	\small
	\textbf{Key} $\kappa_i=m_i-a_i$. \alert{Claim:} an order sorted by decreasing $\kappa$ is optimal.

	\begin{proof}[Exchange argument]
		Suppose $i$ is run immediately before $j$ with $\kappa_i\ge\kappa_j$ versus the swapped order,
		with identical energy before and after the pair. Then
		\begin{align*}
			E(i,j)&=\max\{a_i+\kappa_i,\;a_i+a_j+\kappa_j\}\\
			&\le a_i+a_j+\kappa_i\le \max\{m_j,\;a_j+m_i\}=E(j,i).
		\end{align*}
		Any order with an adjacent pair violating $\kappa_i\ge\kappa_j$ is therefore no better than
		the order with the pair swapped; repeated swaps sort the sequence (bubble sort) without increasing the requirement.
	\end{proof}

	Time $O(n\log n)$. The structure that makes this problem easy is an \alert{exchange argument}, not a state space.
\end{frame}

\begin{frame}{Rapid fire: DP, or not?}
	\audience{Vote}{%
		For each: \textbf{DP}, \textbf{something better}, or \textbf{no compact DP}?
	}
	\small
	\begin{enumerate}
		\item Minimum number of coins summing to $V$.
		\item Unweighted interval scheduling (maximum number of disjoint intervals).
		\item Longest \emph{simple} path in a general directed graph.
		\item $0/1$ knapsack with $n=40$ items and weights up to $10^{12}$.
		\item Longest increasing subsequence.
		\item Shortest paths in a graph with cycles and non-negative weights.
	\end{enumerate}
\end{frame}

\begin{frame}{Rapid fire: answers}
	\small
	\begin{enumerate}
		\item \textbf{DP.} State $=$ remaining amount.
		\item \textbf{Something better.} Earliest-finish-time greedy ($O(n\log n)$); DP is valid but unnecessary.
		\item \textbf{No compact DP.} NP-hard; subset DP is $\Theta(2^nn)$, the state has to record the visited set.
		\item \textbf{Something better.} $O(nW)$ is hopeless; meet-in-the-middle gives $O(2^{n/2}\,n)$.
		\item \textbf{DP, then better.} $O(n^2)$ DP; changing the state to \emph{(length $\to$ smallest tail)}
		gives $O(n\log n)$ by patience sorting.
		\item \textbf{DP in disguise.} The Bellman equations hold but are cyclic; Dijkstra chooses an
		evaluation order. This is the theme of Part~4.
	\end{enumerate}
\end{frame}
